\section{微分形式的积分}
\sourcepage{68--69}
微分1形式可沿曲线积分。若形式闭，则积分仅依赖曲线的同伦类。因此在单连通曲面$X$上，对闭1形式沿起点固定、终点变化的曲线作不定积分，得到$X$上的良定义函数。一般地，闭形式的积分产生多值函数，但其多值性质很特殊，本节将详细研究。我们还考虑2形式的积分，以把线积分转为面积分，并证明留数定理。
\subsection*{A. 微分1形式}
\entry{10.1}{}
设$X$为黎曼曲面，$\omega\in\mathcal E^{(1)}(X)$。设$c:[0,1]\to X$为分段连续可微曲线，即存在分割$0=t_0<\cdots<t_n=1$及坐标卡$(U_k,z_k)$，$z_k=x_k+iy_k$，使$c([t_{k-1},t_k])\subset U_k$，且$x_k\circ c,y_k\circ c$在各子区间上有连续一阶导数。在$U_k$上写$\omega=f_kdx_k+g_kdy_k$，定义
\[
\int_c\omega:=\sum_{k=1}^n\int_{t_{k-1}}^{t_k}
\left(f_k(c(t))\frac{d(x_k\circ c)}{dt}+g_k(c(t))\frac{d(y_k\circ c)}{dt}\right)dt.
\]
容易验证，此定义不依赖分割及坐标卡的选择。
\entry{10.2}{定理}
若$c$分段连续可微，$F\in\mathcal E(X)$，则$\int_c dF=F(c(1))-F(c(0))$。
\begin{proof}
取上述分割与坐标卡。在$U_k$上$dF=F_{x_k}dx_k+F_{y_k}dy_k$，所以
\sourcepage{70}
\begin{align*}
\int_c dF&=\sum_k\int_{t_{k-1}}^{t_k}\left(F_{x_k}(c(t))(x_k\circ c)'(t)+F_{y_k}(c(t))(y_k\circ c)'(t)\right)dt\\
&=\sum_k\int_{t_{k-1}}^{t_k}\frac d{dt}F(c(t))dt
=\sum_k(F(c(t_k))-F(c(t_{k-1})))\\&=F(c(1))-F(c(0)).
\end{align*}
\end{proof}
\entry{10.3}{定义}
若$F\in\mathcal E(X)$满足$dF=\omega$，则称$F$为$\omega$的\textbf{原函数}（primitive）。由(9.15)，有原函数的形式必闭。原函数不唯一：若$F$为原函数，则$F+c$也为原函数。反之，两个原函数只差常数，因为$dF=0$由(10.2)可推出$F$为常数。

由(10.2)，已知原函数便可计算线积分，且恰当微分形式沿曲线的积分只依赖起点、终点。
\entry{10.4}{原函数的局部存在性}
设$U:=\{z:|z|<r\}$，$r>0$，$\omega=f\,dx+g\,dy\in\mathcal E^{(1)}(U)$，并假设$d\omega=0$。由$d\omega=(g_x-f_y)dx\wedge dy$，此条件等价于$g_x=f_y$。我们证明其原函数可取为
\[
F(x,y):=\int_0^1(f(tx,ty)x+g(tx,ty)y)\,dt.
\]
它显然无限次可微，只须验证$F_x=f$、$F_y=g$。积分号下求导得
\[
F_x(x,y)=\int_0^1(tf_x(tx,ty)x+tg_x(tx,ty)y+f(tx,ty))\,dt.
\]
\sourcepage{71}
由$g_x=f_y$及$\frac d{dt}f(tx,ty)=f_x(tx,ty)x+f_y(tx,ty)y$，有
\[
F_x=\int_0^1\left(t\frac d{dt}f(tx,ty)+f(tx,ty)\right)dt
=\int_0^1\frac d{dt}(tf(tx,ty))dt=f(x,y).
\]
同理$F_y=g$，故$dF=\omega$。

若$\omega$全纯，证明更简单：写$\omega=f\,dz$，$f(z)=\sum_{n=0}^\infty c_nz^n$，定义$F(z):=\sum_{n=0}^\infty\frac{c_n}{n+1}z^{n+1}$便得全纯原函数。

整体而言，闭形式的原函数通常仅作为多值函数存在。下一定理使此说法精确化。
\entry{10.5}{定理}
设$\omega\in\mathcal E^{(1)}(X)$闭。则存在连通覆盖$p:\widehat X\to X$及$p^*\omega$的原函数$F\in\mathcal E(\widehat X)$。
\begin{proof}
令$\mathcal F$为$\omega$的原函数层：$\mathcal F(U):=\{f\in\mathcal E(U):df=\omega|_U\}$。由于连通区域上两个原函数只差常数，$\mathcal F$满足恒等定理。由(6.10)，伴随空间$|\mathcal F|$为Hausdorff空间。由(10.4)，每个$a\in X$有连通开邻域$U$及原函数$f$。$f+c$（$c\in\C$）恰为$U$上的所有原函数，因此
\[
p^{-1}(U)=\bigcup_{c\in\C}[U,f+c].
\]
\sourcepage{72}
这些集合两两不交，且投影限制均为到$U$的同胚，故$p:|\mathcal F|\to X$为覆盖。取其任意连通分支$\widehat X$，限制仍为覆盖。$\widehat X$由函数芽组成，故自然定义$F(\varphi):=\varphi(p(\varphi))$，直接由定义可知$dF=p^*\omega$。
\end{proof}
\entry{10.6}{推论}
设$\pi:\widetilde X\to X$为万有覆盖，$\omega$为闭1形式。则$\pi^*\omega$有原函数$f\in\mathcal E(\widetilde X)$。
\begin{proof}
取(10.5)的$p:\widehat X\to X$及原函数$F$。由万有性，存在保纤维全纯映射$\tau:\widetilde X\to\widehat X$。令$f:=\tau^*F$，便是$\tau^*(p^*\omega)=\pi^*\omega$的原函数。
\end{proof}
\entry{10.7}{推论}
单连通黎曼曲面上的每个闭1形式都有原函数。

由(10.6)及$\mathrm{id}:X\to X$为万有覆盖即得。
\entry{10.8}{定理}
设$p:\widetilde X\to X$为万有覆盖，$F$为$p^*\omega$的原函数。若$c:[0,1]\to X$分段连续可微，$\widetilde c$为其提升，则$\int_c\omega=F(\widetilde c(1))-F(\widetilde c(0))$。
\begin{proof}
由定义，对任意分段连续可微曲线$v$有$\int_vp^*\omega=\int_{p\circ v}\omega$，再用(10.2)即得。
\end{proof}
\entry{10.9}{注记}
现在可用(10.8)的公式定义闭形式沿任意连续曲线$c$的积分。它不依赖原函数的选择，因为两个原函数只差常数。它也不依赖提升的选择：若$u,v$为$c$的两个提升，由万有覆盖为Galois覆盖，存在覆盖变换$\sigma$使$v=\sigma\circ u$。
\sourcepage{73}
由$p\circ\sigma=p$，$\sigma^*F$也是$p^*\omega$的原函数，故$\sigma^*F-F$为常数。因此
\[
F(v(1))-F(v(0))=\sigma^*F(u(1))-\sigma^*F(u(0))=F(u(1))-F(u(0)).
\]
\entry{10.10}{定理}
设$\omega$为闭1形式。
(a) 若$u,v$为从$a$到$b$的同伦曲线，则$\int_u\omega=\int_v\omega$；
(b) 若$u,v$为自由同伦的闭曲线，则也有此等式。
\begin{proof}
(a) 在万有覆盖中取同起点的提升$\widetilde u,\widetilde v$。由(4.10)，终点亦同，由(10.8)即得。

(b) 若$u,v$的基点分别为$x_0,x_1$，由(3.13)存在从$x_0$到$x_1$的曲线$w$使$u\sim w\cdot v\cdot w^-$。由(a)，
\[
\int_u\omega=\int_{w\cdot v\cdot w^-}\omega=\int_w\omega+\int_v\omega-\int_w\omega=\int_v\omega.
\]
\end{proof}
\entry{10.11}{周期}
由(10.10)，对$\sigma\in\pi_1(X)$可以取任意代表曲线$u$定义$a_\sigma:=\int_\sigma\omega:=\int_u\omega$。这些积分称为$\omega$的\textbf{周期}（period）。显然$\int_{\sigma\tau}\omega=\int_\sigma\omega+\int_\tau\omega$，因而得到群同态$\pi_1(X)\to(\C,+)$，称为$\omega$的\textbf{周期同态}。
\sourcepage{74}
\textbf{例。}$X=\C^*$时，$\pi_1(X)\cong\Z$由$u(t)=e^{2\pi it}$的类生成。对$\omega=dz/z$有$\int_u\omega=2\pi i$，故周期同态为$\Z\to\C$，$n\mapsto2\pi in$；这里同构$\Z\to\pi_1(\C^*)$取$n\mapsto\operatorname{cl}(u^n)$。
\entry{10.12}{自守加项}
设$p:\widetilde X\to X$为万有覆盖，$G:=\operatorname{Deck}(\widetilde X/X)\cong\pi_1(X)$。对$\sigma\in G$及函数$f:\widetilde X\to\C$，令$\sigma f:=f\circ\sigma^{-1}$。则$\sigma(f+g)=\sigma f+\sigma g$、$\sigma(fg)=(\sigma f)(\sigma g)$及$(\sigma\tau)f=\sigma(\tau f)$。

若存在常数$a_\sigma\in\C$使$f-\sigma f=a_\sigma$，则称$f$\textbf{具有常自守加项的加法自守函数}，这些唯一确定的常数称为$f$的\textbf{自守加项}（summands of automorphy）。由$f-\tau f=a_\tau$得$\sigma f-\sigma\tau f=a_\tau$，所以
\[
a_{\sigma\tau}=(f-\sigma f)+(\sigma f-\sigma\tau f)=a_\sigma+a_\tau.
\]
因此$\sigma\mapsto a_\sigma$是$G\to\C$的同态。

覆盖变换不变函数（对每个$\sigma$有$\sigma f=f$）是自守加项全为零的例子。这时存在$f_0:X\to\C$使$f=p^*f_0$；若$f$可微或全纯，$f_0$也相应可微或全纯。
\entry{10.13}{定理}
(a) 若$\omega$为闭1形式，$F$为$p^*\omega$的原函数，则$F$具有常自守加项。在同构$\pi_1(X)\cong\operatorname{Deck}(\widetilde X/X)$下，这些加项恰为$\omega$的周期。
\sourcepage{75}
(b) 反之，若$F\in\mathcal E(\widetilde X)$具有常自守加项，则存在唯一闭1形式$\omega$使$dF=p^*\omega$。
\begin{proof}
(a) 由$p\circ\sigma^{-1}=p$，$\sigma F$也是$p^*\omega$的原函数，故$-a_\sigma:=\sigma F-F$为常数。固定$p(z_0)=x_0$，取从$y_0:=\sigma^{-1}(z_0)$到$z_0=\sigma(y_0)$的曲线$v$，则$u:=p\circ v$代表与$\sigma$对应的基本群元素。由(10.8)，
\[
\int_u\omega=F(v(1))-F(v(0))=F(z_0)-F(\sigma^{-1}(z_0))=a_\sigma.
\]
\translatornote{原文最后一项印为$-a_\sigma$，与(10.12)及本段定义相反；这里按前式计算写为$a_\sigma$。}
(b) 由$F$的自守性质，$\sigma^*(dF)=d(\sigma^*F)=d(F+a_\sigma)=dF$。故$dF$在覆盖变换下不变。由$p$局部双全纯，存在唯一$\omega\in\mathcal E^{(1)}(X)$使$p^*\omega=dF$，且它显然闭。
\end{proof}
\entry{10.14}{例}
设$\Gamma=\Z\gamma_1+\Z\gamma_2$为格，$X=\C/\Gamma$。典范投影$\pi:\C\to X$为万有覆盖，覆盖变换为按$\gamma\in\Gamma$的平移。恒等函数$z:\C\to\C$具有自守加项$a_\gamma=\gamma$，故$dz$在覆盖变换下不变。因此存在$\omega\in\Omega(X)$使$\pi^*\omega=dz$，其周期恰为格$\Gamma$的元素。
\entry{10.15}{定理}
闭1形式$\omega$有整体原函数，当且仅当所有周期均为零。
\begin{proof}
若有原函数，(10.2)表明周期全零。反之，由(10.6)在万有覆盖上有$p^*\omega$的原函数$F$。它的自守加项全零，\translatornote{原文此处引用(10.3)，相关自守加项结论应参见(10.13)。}故存在$f\in\mathcal E(X)$使$F=p^*f$。
\sourcepage{76}
由$p^*\omega=dF=d(p^*f)=p^*(df)$得$\omega=df$。
\end{proof}
\textbf{注记。}周期全零时，可固定任意$x_0\in X$，取$f(x):=\int_{x_0}^x\omega$为原函数；积分沿任意连接曲线，与曲线选择无关。
\entry{10.16}{推论}
若紧黎曼曲面上的两个全纯微分形式$\omega_1,\omega_2$具有相同周期同态，则$\omega_1=\omega_2$。
\begin{proof}
$\omega:=\omega_1-\omega_2$周期全零，故有全纯原函数$f$。由紧性$f$为常数，故$\omega=df=0$。
\end{proof}
\subsection*{B. 微分2形式}
\entry{10.17}{}
先考虑复平面上的积分。设$U\subset\C$开，$\omega\in\mathcal E^{(2)}(U)$，写$\omega=f\,dx\wedge dy=\frac i2f\,dz\wedge d\bar z$。若$f$在某个紧子集之外为零，定义$\iint_U\omega:=\iint_Uf(x,y)\,dx\,dy$，右端为通常二重积分。

若$\varphi:V\to U$双全纯，写$\varphi=u+iv$，由Cauchy--Riemann方程，其Jacobi行列式$\partial(u,v)/\partial(x,y)=u_xv_y-u_yv_x=|\varphi'|^2$。故变量代换公式为
\[
\iint_Uf\,dx\,dy=\iint_V(f\circ\varphi)|\varphi'|^2dx\,dy.
\]
\sourcepage{77}
另一方面$\varphi^*(dz\wedge d\bar z)=d\varphi\wedge d\bar\varphi=|\varphi'|^2dz\wedge d\bar z$，因此$\varphi^*\omega=(f\circ\varphi)|\varphi'|^2dx\wedge dy$，从而$\iint_U\omega=\iint_V\varphi^*\omega$。
\entry{10.18}{}
在黎曼曲面$X$上，微分形式$\omega$的\textbf{支集}定义为闭集$\operatorname{Supp}(\omega):=\overline{\{a\in X:\omega(a)\ne0\}}$。函数支集类似定义。

(a) 若坐标卡$\varphi:U\to V$包含$\omega$的紧支集，则$(\varphi^{-1})^*\omega$在$V\subset\C$中有紧支集，定义
\[
\iint_X\omega:=\iint_U\omega:=\iint_V(\varphi^{-1})^*\omega.
\]
它与坐标无关：若另一坐标卡$\varphi_1:U_1\to V_1$也包含支集，取交集可假设$U=U_1$。则$\psi:=\varphi_1\circ\varphi^{-1}:V\to V_1$双全纯，且$(\varphi^{-1})^*\omega=\psi^*((\varphi_1^{-1})^*\omega)$，由(10.17)两个积分相等。

(b) 若$\omega$为任意紧支集2形式，用有限张坐标卡$\varphi_k:U_k\to V_k$覆盖支集。可取\textbf{单位分解}（附录A）$f_k\in\mathcal E(X)$满足$\operatorname{Supp}(f_k)\subset U_k$，且$\sum_kf_k(x)=1$对每个$x\in\operatorname{Supp}(\omega)$成立。
\sourcepage{78}
则$\omega=\sum_kf_k\omega$，每项有包含于$U_k$的紧支集。定义$\iint_X\omega:=\sum_k\iint_Xf_k\omega$。右端按(a)定义，容易验证与坐标卡及单位分解的选择无关。
\entry{10.19}{}
后面需要平面Stokes定理的特殊情形：若$U\subset\C$开，$A\subset U$紧且有光滑边界$\partial A$，则对每个$\omega\in\mathcal E^{(1)}(U)$，有$\iint_A d\omega=\int_{\partial A}\omega$。边界定向约定为：先取指向$A$外部的法向量，再取边界切向量，构成平面正定向基。

我们只需圆盘或圆环$A=\{z:\varepsilon\le|z|\le R\}$（$0<\varepsilon<R$）的情形。圆环边界由正向外圆$|z|=R$和负向内圆$|z|=\varepsilon$组成。若$\omega=f\,dx+g\,dy$，则定理为
\[
\iint_A(g_x-f_y)dx\,dy=\int_{|z|=R}(f\,dx+g\,dy)-\int_{|z|=\varepsilon}(f\,dx+g\,dy).
\]
直接引入极坐标$x=r\cos\theta$、$y=r\sin\theta$证明。先令$\omega=g\,dy$，则$d\omega=g_xdx\wedge dy$，并有$\partial_x=\cos\theta\,\partial_r-(\sin\theta/r)\partial_\theta$。
\sourcepage{79}
令$\widetilde g(r,\theta):=g(re^{i\theta})$，则
\begin{align*}
\iint_A d\omega&=\int_\varepsilon^R\int_0^{2\pi}(\cos\theta\,\widetilde g_r-\tfrac{\sin\theta}r\widetilde g_\theta)r\,d\theta\,dr\\
&=\int_\varepsilon^R\int_0^{2\pi}\left(\cos\theta\,\partial_r(r\widetilde g)-\partial_\theta(\sin\theta\widetilde g)\right)d\theta\,dr.
\end{align*}
对固定$r$，第二项的$\theta$积分为$[\sin\theta\widetilde g]_0^{2\pi}=0$。所以
\begin{align*}
\iint_A d\omega&=\int_0^{2\pi}\cos\theta\left(\int_\varepsilon^R\partial_r(r\widetilde g)dr\right)d\theta\\
&=\int_0^{2\pi}\widetilde g(R,\theta)R\cos\theta d\theta-
\int_0^{2\pi}\widetilde g(\varepsilon,\theta)\varepsilon\cos\theta d\theta\\
&=\int_{|z|=R}g\,dy-\int_{|z|=\varepsilon}g\,dy=\int_{\partial A}\omega.
\end{align*}
$\omega=f\,dx$的情形经坐标变换$(x,y)\mapsto(y,-x)$化为上例，此变换Jacobi行列式为1。故圆环情形得证，令$\varepsilon\to0$得到圆盘情形。
\entry{10.20}{定理}
若$\omega\in\mathcal E^{(1)}(X)$有紧支集，则$\iint_Xd\omega=0$。
\begin{proof}
用单位分解可写$\omega=\omega_1+\cdots+\omega_n$，每项的紧支集包含于一张坐标卡。故可假设$X=\C$。取$R$足够大使$\operatorname{Supp}(\omega)\subset\{|z|<R\}$，则
\[
\iint_\C d\omega=\iint_{|z|\le R}d\omega=\int_{|z|=R}\omega=\int_{|z|=R}0=0.
\]
\end{proof}
\sourcepage{80}
\entry{10.21}{留数定理}
设$X$为紧黎曼曲面，$a_1,\ldots,a_n$互异，$X':=X\setminus\{a_1,\ldots,a_n\}$。则对每个$\omega\in\Omega(X')$，有$\sum_{k=1}^n\operatorname{Res}_{a_k}\omega=0$。
\begin{proof}
选两两不交坐标邻域$(U_k,z_k)$，使$z_k(a_k)=0$且$z_k(U_k)$为圆盘。取紧支集函数$f_k$，$\operatorname{Supp}(f_k)\subset U_k$，并在$a_k$的某开邻域$V_k\subset U_k$上恒为1。令$g:=1-\sum_kf_k$，则$g|_{V_k}=0$。把$g\omega$在各$a_k$处延拓为零，得到$\mathcal E^{(1)}(X)$中元素。由(10.20)，$\iint_Xd(g\omega)=0$。

由于$d\omega=0$，在$V_k\cap X'$上$d(f_k\omega)=d\omega=0$，故$d(f_k\omega)$可视为$X$上的2形式，其支集紧包含于$U_k\setminus\{a_k\}$。由$d(g\omega)=-\sum_kd(f_k\omega)$，得$\sum_k\iint_Xd(f_k\omega)=0$。

只须证明$\iint_Xd(f_k\omega)=-2\pi i\operatorname{Res}_{a_k}\omega$。积分只需在$U_k$进行。利用$z_k$把$U_k$等同于单位圆盘，取$0<\varepsilon<R<1$，使$f_k=1$于$|z_k|\le\varepsilon$，$f_k=0$于$|z_k|\ge R$。由Stokes定理及平面留数定理，
\[
\iint_Xd(f_k\omega)=\iint_{\varepsilon\le|z_k|\le R}d(f_k\omega)
=\int_{|z_k|=R}f_k\omega-\int_{|z_k|=\varepsilon}f_k\omega
=-\int_{|z_k|=\varepsilon}\omega=-2\pi i\operatorname{Res}_{a_k}\omega.
\]
\end{proof}
\entry{10.22}{推论}
紧黎曼曲面上的非常值亚纯函数，按重数计算，零点数等于极点数。
\sourcepage{81}
\begin{proof}
$\omega:=df/f$除$f$的零点、极点外均全纯。在$m$阶零点处留数为$m$，在$m$阶极点处为$-m$。由留数定理即得。
\end{proof}
\textbf{注记。}(4.25)已用覆盖证明过此结论。
\subsection*{习题（第10节）}
\entry{10.1}{}
设$\omega$为$X$上的全纯1形式，$\varphi$为$a\in X$某邻域上的原函数，$(Y,p,f,b)$为其极大解析延拓。证明：(a)$p$为覆盖；(b)$f$为$p^*\omega$的原函数；(c)$p$为Galois覆盖且$\operatorname{Deck}(Y/X)$为阿贝尔群。
\entry{10.2}{}
设$X=\C/\Gamma$为环面。对任意同态$a:\pi_1(X)\to\C$，证明存在闭1形式$\omega\in\mathcal E^{(1)}(X)$使周期同态为$a$。
\entry{10.3}{}
设$\omega\in\MM^{(1)}(X)$在每个极点处留数为零。证明存在覆盖$p:Y\to X$及亚纯函数$F\in\MM(X)$使$dF=p^*\omega$。\translatornote{原文写$F\in\MM(X)$；由等式的定义域，应为$F\in\MM(Y)$。}
\entry{10.4}{}
设$\Gamma\subset\C$为格。用留数定理证明，不存在恰有一个一阶极点的亚纯函数$f\in\MM(\C/\Gamma)$。
